\chapter{把生成模型放回科学计算}
\label{chap:science}

\begin{keybox}[把分布问题换回物理问题]
前面各章的模型都可以放回科学计算：能量给出目标分布，样本给出构型，密度给出自由能。这一章看几个已经走通的接口——有限温采样、重整化群、分子与晶体结构、量子线路与量子态——并在每个接口上说清楚：模型生成的对象是什么，评价它的物理量又是什么。
\end{keybox}

\section{统计物理：从能量到平衡样本}

以 Ising 自旋为例，一个构型 \(\sigma=(\sigma_1,\ldots,\sigma_N)\) 的能量由相互作用和外场给定，有限温的目标分布是 \(p_\beta(\sigma)=e^{-\beta H(\sigma)}/Z_\beta\)，其中 \(\beta=1/(k_{\mathrm B}T)\)。难点是构型数随系统大小指数增长，配分函数 \(Z_\beta\) 和观测量的平衡平均通常无法直接求和。生成模型可以构造容易采样的试探分布，再用已知能量引导它接近目标分布。

对连续的多体坐标，Boltzmann generator 用可逆流把易采样的基分布映射到目标构型空间~\cite{kim2019boltzmann}。模型给出候选样本及其密度后，可以按目标玻尔兹曼权重重新加权，或用 MCMC 修正偏差。模型的直接作用是快速产生物理上合理的候选构型；加权后的样本是否代表目标温度下的平衡测度，是另一个问题。

若用变分自由能训练 Boltzmann generator，抽样从参数无关的基分布开始，连续构型由可逆流映射得到；能量和模型对数密度的梯度沿这一映射回传，使用第 \ref{sec:gradient-estimator-map} 节的重参数化路径。

生成模型还可以当重要性采样的 proposal。设从 \(q_\theta\) 采样 \(x\)，目标权重是 \(w(x)=e^{-\beta E(x)}/q_\theta(x)\)；加权平均给出目标测度下的观测量，配分函数由权重的均值估计。若 \(q_\theta\) 只在部分区域有概率，权重的方差会很大；退火重要性采样把一列中间温度串起来，让相邻分布的重叠更大，得到更稳的自由能差估计。这里生成模型提供覆盖多个模式的 proposal，最终的自由能仍由权重和物理能量决定。

\section{有限温变分计算：自回归网络}

传统变分平均场常把试探分布写成独立单点概率的乘积，因而难以描述自旋之间的关联。变分自回归网络（Variational Autoregressive Network，VAN）保留链式分解 \(q_\theta(\sigma)=\prod_i q_\theta(\sigma_i\mid\sigma_{<i})\)，但让后面的条件概率读取前面已生成的自旋~\cite{wu2019variational}。每个条件概率都归一化，所以完整构型的 \(q_\theta(\sigma)\) 可计算；逐个抽取自旋便得到独立生成的完整构型，无需等待马尔可夫链混合。

已知 \(H\) 时，可以直接最小化变分自由能 \(F_\theta=\mathbb E_{q_\theta}[H(\sigma)+\beta^{-1}\log q_\theta(\sigma)]\)。它满足 \(F_\theta-F_\beta=\beta^{-1}\KL(q_\theta\Vert p_\beta)\geq 0\)，因此给出真实自由能的上界，也把降低能量与保留熵放在同一目标中。训练后，用 \(\sigma\sim q_\theta\) 的样本估计能量、磁化和关联函数；若要准确估计目标玻尔兹曼平均，还需考虑分布误差并使用适当的重加权或采样修正。VAN 是比独立变量平均场更灵活的变分分布，优化目标仍是有限温的自由能。

VAN 从依赖 \(\theta\) 的离散分布 \(q_\theta(\sigma)\) 抽取自旋，无法沿自旋取值反向传播；训练自由能时使用第 \ref{sec:gradient-estimator-map} 节的 score-function 路径，以构型的局部自由能 \(H(\sigma)+\beta^{-1}\log q_\theta(\sigma)\) 加权 \(\nabla_\theta\log q_\theta(\sigma)\)，并可减去 baseline 降低方差。

\section{流模型与重整化群}

重整化群试图逐层分离短尺度涨落，留下描述长尺度行为的有效变量。Li 与 Wang 的神经网络重整化群用分层可逆变换连接物理构型和多尺度潜变量~\cite{li2018neuralrg}：每一层提取部分自由度，并让潜变量之间的统计依赖逐渐减弱；反向变换则从简单的独立潜变量生成物理构型。这使“粗粒化后保留什么信息”成为可以通过模型学习的问题。

可逆性还允许通过 Jacobian 计算生成构型的概率密度，并用变分自由能训练该映射。训练后的潜变量按层对应不同尺度的涨落，可以检查它们的相关性，观察模型是否实现了预期的尺度分离。流模型在这里承担可计算密度与生成样本的工具；分层结构和尺度分离的检验，才赋予这一具体模型重整化群的解释。

对这类连续可逆流，变分自由能中的模型样本由固定基分布经可微映射生成，训练仍沿第 \ref{sec:gradient-estimator-map} 节的重参数化路径求导。

\section{分子构象与材料结构}

若已知分子的原子种类和连接关系，生成对象可以是所有原子的三维坐标；扩散模型或 Flow Matching 从简单噪声逐步得到候选构象。模型应尊重平移、旋转及同种原子置换带来的等价性，同时避免不合理的键长、碰撞和错误手性。等变架构把群作用写进网络结构，例如对坐标做等变的消息传递，或在损失中加入不变量；约束可以用参数化保证、投影、惩罚或后验筛选实现。评估不能只看坐标是否像训练样本，还要检查能量、构象多样性及目标性质，并说明这些量是在生成后直接计算还是经过物理松弛得到。

晶体生成还需要处理元素组成、周期性晶胞和晶胞内原子位置。若给定目标性质作为条件，得到的结构仍须满足组成、对称性和稳定性约束，再由独立的物理计算验证性质。这个例子说明，条件生成给出的是候选分布，后续筛选与验证决定候选能否成为可信的科学结果。

\section{可微量子线路结构搜索}

设计变分量子线路时，既要选择每一位置放什么门，也要调节门的连续参数。门的种类和连接方式是离散变量，直接枚举所有线路会迅速遇到组合爆炸。可微量子架构搜索（Differentiable Quantum Architecture Search，DQAS）先限定门集合、线路位置和可用连接，再用可训练的概率分布 \(p_\alpha(A)\) 抽样线路结构 \(A\)~\cite{zhang2022dqas}。对抽到的线路计算任务损失 \(L(A,\theta)\)，例如制备态的能量或逼近目标酉矩阵的误差；优化 \(\mathbb E_{A\sim p_\alpha}[L(A,\theta)]\) 中的结构分布参数 \(\alpha\) 与门参数 \(\theta\)，最后选出具体线路并检验其性能。

DQAS 的“可微”是指期望损失对概率分布参数可以求梯度；一次离散的门选择本身没有连续导数。概率模型让搜索在门组合之间分配计算预算，并把线路搭建与参数优化联系起来。得到的线路仍需在实际任务上重新评估，因为搜索空间限制和抽样梯度的噪声都会影响结果。

优化线路结构分布 \(p_\alpha(A)\) 时，抽到的门结构 \(A\) 是离散的，因此其参数 \(\alpha\) 使用第 \ref{sec:gradient-estimator-map} 节的 score-function 路径，以 \(L(A,\theta)\nabla_\alpha\log p_\alpha(A)\) 估计梯度；给定线路后的连续门参数 \(\theta\) 则通过相应的量子线路梯度方法更新。

\section{神经网络量子态：基态与动力学}

前一节的模型生成线路结构；如果目标是多体系统的量子态本身，可以用神经网络表示波函数振幅 \(\psi_\theta(s)=\langle s\mid\psi_\theta\rangle\)，其中 \(s\) 是选定基底中的构型~\cite{carleo2017nqs}。求基态时，按 \(|\psi_\theta(s)|^2\) 采样构型，用哈密顿量计算变分能量 \(E_\theta=\langle\psi_\theta|H|\psi_\theta\rangle/\langle\psi_\theta|\psi_\theta\rangle\)，再调节网络参数降低能量。变分原理保证 \(E_\theta\) 不低于真实基态能量；能量之外还应检查相关函数等目标观测量。

研究动力学时，从初态出发，把薛定谔方程给出的波函数变化投影到网络参数能够表示的变化方向，再按时间步更新网络参数~\cite{carleo2017nqs}。这种时间依赖变分方法用同一波函数表示跟踪实时演化，但长期精度取决于网络表达能力、采样误差和时间积分误差。这里生成模型的角色是表示带相位的量子振幅；它与只生成某个测量基下比特串概率的 Born machine 有不同的物理目标。

对按 \(|\psi_\theta(s)|^2\) 抽到的离散构型，变分蒙特卡洛也用对数导数把采样分布对参数的依赖写入能量梯度；这与第 \ref{sec:gradient-estimator-map} 节的 score-function 路径相通，权重由哈密顿量给出的局部能量决定。

\section{量子与张量网络生成模型}

前一节用波函数研究量子态本身；若只想生成选定基底下的构型，可以把振幅的模方当作概率分布。同一 Born 规则可以配两种表示：参数化量子线路在测量时直接给出样本，矩阵乘积态（Matrix Product State，MPS）则用一串可收缩的张量表示振幅。两者都可以作为离散构型的生成模型，但训练和采样所需的计算不同。

\textbf{量子线路 Born machine。}QCBM~\cite{liu2018qcbm} 把 \(U_\theta|0^n\rangle\) 测得比特串 \(s\) 的概率写为 \(p_\theta(s)=|\langle s|U_\theta|0^n\rangle|^2\)。第 \ref{sec:mmd-qcbm} 节已说明：运行线路即可采样，真实样本与测量样本可用 MMD 比较，参数移位法则用多次线路运行估计训练梯度。放到多自旋构型任务时，除分布差异外，还应检查磁化和关联函数等任务相关观测量。有限测量次数难以准确估计稀有构型的逐点概率，这是用样本差异训练时仍需注意的限制。

\textbf{张量网络 Born machine。}对 \(s=(s_1,\ldots,s_n)\)，MPS 把振幅写为一串矩阵的乘积，再由模方得到归一化概率~\cite{han2018mps}：
\begin{equation}
 \psi_\theta(s)=A^{[1]}_{s_1}A^{[2]}_{s_2}\cdots A^{[n]}_{s_n},\qquad
 p_\theta(s)=\frac{|\psi_\theta(s)|^2}{\sum_{s'}|\psi_\theta(s')|^2}.
\end{equation}
这里首尾张量的边界维度为 1，所以矩阵乘积最终是一个标量；相邻张量之间的内部维度称为键维 \(D\)。给定构型时，可直接收缩矩阵链计算其概率；生成时则逐位计算条件概率 \(p_\theta(s_i\mid s_{<i})\)、采样并继续收缩，直到得到整条构型。训练可最小化数据的负对数似然，梯度沿张量收缩回传。较大的 \(D\) 容许更多跨位置关联，也增加存储和收缩成本；固定的小键维适合主要由较短程或较低纠缠结构决定的分布。与量子线路的测量采样相比，MPS 的逐点概率和归一化可以通过经典张量收缩计算，但成本仍随键维增长。

这两类 Born machine 学到的都是所选测量基下的 \(p_\theta(s)\)。相同的概率可以对应不同相位的波函数，因此只检验该基底的样本统计，不能判断完整量子态是否正确；若目标是量子态，还需额外测量基或对相位敏感的观测量。

回到第 \ref{chap:probability-language} 章的问题：给定目标性质 \(y^\star\)，怎样在合理的结构上采样 \(p(x\mid y^\star)\)。这一章的例子给出了同一条链的不同实现方式：生成模型提供 \(p(x)\) 或 \(p(x\mid c)\)，性质预测器与物理计算提供 \(p(y\mid x)\) 的近似或验证，权重、引导、筛选和重加权把两者接起来。链条的每一环都可能引入误差，因此最终报告的应是经过独立物理检验的候选。哪一类模型适合哪个接口，取决于对象是离散还是连续、能量是否可算、以及能否承受 MCMC 或 ODE 的采样成本。
